\pagestyle{empty}
\tight
\begin{tblr}{width=\textwidth,colspec={|Q[c,m,1.9cm]|X[l,m]|},hlines,vlines,hspan=minimal,row{1}={bg=headblue},rowsep=1pt,colsep=3pt}
\SetCell{c,m}\hfil\logo\hfil & \textbf{POLITEKNIK NEGERI MALANG}\par\textbf{JURUSAN TEKNOLOGI INFORMASI}\par\textbf{PROGRAM STUDI: D4 TEKNIK INFORMATIKA} \\
\SetCell[c=2]{c} \textbf{RENCANA PEMBELAJARAN SEMESTER (RPS)} & \\
\end{tblr}
\begin{longtblr}{width=\textwidth,colspec={|Q[l,m,1.9cm]|X[0.85,l,m]|X[1.45,l,m]|X[0.75,l,m]|X[1.15,l,m]|},hlines,vlines,hspan=minimal,rowsep=1pt,colsep=2pt,row{1,3}={bg=headgray,font=\bfseries}}
MATA KULIAH & KODE & BOBOT (SKS)/jam & SEMESTER & TGL. PENYUSUNAN \\
Desain Antarmuka & RTI252003 & 2 SKS/4 jam & 2 & 29 Januari 2026 \\
OTORISASI & Dosen Pengembang RPS & Koordinator RMK & \SetCell[c=2]{l} Ka PRODI &  \\
 & Anugrah Nur Rahmanto, S.Sn., M.Ds. & Anugrah Nur Rahmanto, S.Sn., M.Ds. & \SetCell[c=2]{l}{\vspace{8mm}\par Dr. Ely Setyo Astuti, S.T., M.T.} &  \\
\SetCell[r=13]{l} Capaian Pembelajaran (CP) & \SetCell[c=4]{l,bg=headgray,font=\bfseries} Capaian Pembelajaran Lulusan Program Studi (CPL-Prodi) &  &  &  \\
 & CPL02 & \SetCell[c=3]{l} Mampu menerapkan prinsip rekayasa sistem dan perangkat lunak untuk menghasilkan solusi aplikasi multi-platform sesuai kebutuhan. &  &  \\
 & CPL06 & \SetCell[c=3]{l} Mampu merancang, mengimplementasikan, menguji, melakukan deployment, memelihara, serta menjamin mutu perangkat lunak yang menjawab permasalahan dan memenuhi kebutuhan pengguna. &  &  \\
 & CPL08 & \SetCell[c=3]{l} Mampu mengelola proyek teknologi informasi secara profesional melalui perencanaan, koordinasi, komunikasi, dan optimalisasi sumber daya. &  &  \\
 & \SetCell[c=4]{l,bg=headgray,font=\bfseries} Capaian Pembelajaran mata kuliah (CPL-MK) &  &  &  \\
 & CPMK0202 & \SetCell[c=3]{l} Mahasiswa mampu menerapkan prinsip dan metode UX Design untuk merancang prototipe antarmuka aplikasi multi-platform yang interaktif. &  &  \\
 & CPMK0601 & \SetCell[c=3]{l} Mahasiswa mampu merancang dan mengevaluasi kualitas pengalaman pengguna melalui prototipe yang menjawab masalah interaksi pengguna. &  &  \\
 & CPMK0802 & \SetCell[c=3]{l} Mahasiswa mampu mengelola proses rekayasa perangkat lunak dan UI/UX dalam siklus hidup proyek untuk menghasilkan produk berpusat pada pengguna. &  &  \\
 & \SetCell[c=4]{l,bg=headgray,font=\bfseries} Kemampuan akhir tiap tahapan belajar (Sub-CPMK) &  &  &  \\
 & SCPMK0202-01201 & \SetCell[c=3]{l} Mahasiswa mampu menjelaskan prinsip UI/UX, design thinking, user persona, dan journey map untuk mengidentifikasi kebutuhan pengguna. &  &  \\
 & SCPMK0202-01202 & \SetCell[c=3]{l} Mahasiswa mampu merancang wireframe, visual design, design system, dan UX writing berdasarkan kebutuhan pengguna. &  &  \\
 & SCPMK0601-01203 & \SetCell[c=3]{l} Mahasiswa mampu menghasilkan prototipe interaktif dan mengevaluasi usability berdasarkan skenario pengujian. &  &  \\
 & SCPMK0802-01204 & \SetCell[c=3]{l} Mahasiswa mampu mengelola artefak proyek UI/UX, mempresentasikan rationale desain, dan menyusun rekomendasi perbaikan. &  &  \\
\SetCell[r=2]{l} Penilaian & \SetCell[c=4]{l} *Nilai CPMK didapat dari agregasi nilai SUB-CPMK yang dibebankan &  &  &  \\
 & \SetCell[c=4]{l}{\tiny\begin{tblr}{width=\linewidth,colspec={|Q[c,m,0.1\linewidth]|Q[c,m,0.16\linewidth]|X[l,m]|X[l,m]|X[l,m]|X[l,m]|X[l,m]|X[l,m]|Q[c,m,0.08\linewidth]|},hlines,vlines,hspan=minimal,rowsep=1pt,colsep=0.7pt,row{1,2}={bg=headgray,font=\bfseries}}
Kode CPMK & Kode SCPMK & \SetCell[c=6]{c} Bobot per Bentuk Penilaian & & & & & & Total Bobot \\
 & & Tugas & Kuis 1 & UTS & Kuis 2 & PBL & UAS & \\
\SetCell[r=2]{c} CPMK0202 & SCPMK0202-01201 & 8.40\%: UI/UX, design thinking, persona & 5\%: konsep UI/UX & 12.85\%: artefak desain awal & & & & 26.25\% \\
 & SCPMK0202-01202 & 14.00\%: Figma, wireframe, visual, design system, UX writing & & 7.15\%: artefak desain awal & & 5.70\%: proyek UI/UX & & 26.85\% \\
\SetCell[r=1]{c} CPMK0601 & SCPMK0601-01203 & 11.20\%: prototype, UX motion, usability testing & & & 5\%: design system dan prototype & & 12.55\%: validasi prototype & 28.75\% \\
\SetCell[r=1]{c} CPMK0802 & SCPMK0802-01204 & 2.80\%: iterasi desain & & & & 5.70\%: proyek UI/UX & 9.65\%: validasi prototype & 18.15\% \\
\SetCell[c=8]{r}\textbf{Total Per Penilaian} & & & & & & & & \textbf{100\%} \\
\end{tblr}} &  &  &  \\
Diskripsi Singkat Mata Kuliah & \SetCell[c=4]{l} Mata kuliah Desain Antarmuka membangun kemampuan merancang UI/UX melalui riset pengguna, design thinking, persona, journey map, wireframe, visual design, design system, prototipe interaktif, motion, dan usability testing. &  &  &  \\
Bahan Kajian / Materi Pembelajaran & \SetCell[c=4]{l}\begin{enumerate}\item Pengenalan UI/UX dan design thinking\item Persona, journey map, dan kebutuhan pengguna\item Figma, wireframe, visual design, dan design system\item UX writing, prototyping, UX motion, dan usability testing\end{enumerate} &  &  &  \\
\SetCell[r=2]{l} Pustaka & \SetCell[c=4]{l} \textbf{Utama:}\begin{enumerate}\item Buley, L. 2013. The User Experience Team of One. Rosenfeld Media.\item Young, I. 2008. Mental Models. Rosenfeld Media.\item Frost, B. 2016. Atomic Design.\end{enumerate} &  &  &  \\
 & \SetCell[c=4]{l} \textbf{Pendukung:} Materi LMS, dokumentasi resmi perangkat lunak, dan jobsheet praktikum. &  &  &  \\
Media Pembelajaran & \SetCell[c=2]{l} \textbf{Software:} Figma, browser, office suite, LMS. &  & \SetCell[c=2]{l} \textbf{Hardware:} Komputer/laptop, proyektor. &  \\
Nama Dosen Pengampu & \SetCell[c=4]{l}\begin{enumerate}\item Anugrah Nur Rahmanto, S.Sn., M.Ds.\item Farid Angga Pribadi, S.Kom., M.Kom.\item Bagas Satya Dian Nugraha, S.T., M.T.\item Dimas Wahyu Wibowo, S.T., M.T.\item Habibie Ed Dien, S.Kom., M.T.\end{enumerate} &  &  &  \\
Matakuliah Syarat & \SetCell[c=4]{l} - &  &  &  \\
\end{longtblr}
\begin{longtblr}{width=\textwidth,colspec={|Q[c,m,0.06\textwidth]|X[1.35,l,m]|X[1.08,l,m]|X[1.16,l,m]|X[0.7,l,m]|X[1.24,l,m]|X[1.06,l,m]|X[0.98,l,m]|Q[c,m,0.05\textwidth]|},hlines,vlines,hspan=minimal,rowhead=2,row{1,2}={bg=headgreen,font=\bfseries},rowsep=1pt,colsep=1.5pt}
Mgg.\par Ke & Kemampuan Akhir Yang Direncanakan (Sub-CP-MK) & Bahan kajian (Materi Pembelajaran) & Bentuk dan Metode Pembelajaran & Estimasi Waktu & Pengalaman Belajar Mahasiswa & Kriteria \& Bentuk Penilaian & Indikator Penilaian & Bobot Penilaian (\%) \\
(1) & (2) & (3) & (4) & (5) & (6) & (7) & (8) & (9) \\
1 & SCPMK0202-01201 & Pengenalan UI/UX. & Modalitas: Blended Learning. Bentuk: Luring/Daring. Metode: Design Thinking, studio review, studi kasus, dan peer feedback. & 1 x 2 x 50' tatap muka; 1 x 2 x 50' tugas/praktik mandiri. & Mahasiswa mengerjakan aktivitas terarah untuk pengenalan ui/ux dan menunjukkan evidence hasil belajar. & Bentuk penilaian: Pengenalan UI/UX. Mengacu rubrik RTM. & Ketepatan konsep; kualitas artefak/praktik; validasi dan dokumentasi. & 2.8 \\
2 & SCPMK0202-01201 & Design thinking. & Modalitas: Blended Learning. Bentuk: Luring/Daring. Metode: Design Thinking, studio review, studi kasus, dan peer feedback. & 1 x 2 x 50' tatap muka; 1 x 2 x 50' tugas/praktik mandiri. & Mahasiswa mengerjakan aktivitas terarah untuk design thinking dan menunjukkan evidence hasil belajar. & Bentuk penilaian: Design thinking. Mengacu rubrik RTM. & Ketepatan konsep; kualitas artefak/praktik; validasi dan dokumentasi. & 2.8 \\
3 & SCPMK0202-01201 & Persona dan journey map. & Modalitas: Blended Learning. Bentuk: Luring/Daring. Metode: Design Thinking, studio review, studi kasus, dan peer feedback. & 1 x 2 x 50' tatap muka; 1 x 2 x 50' tugas/praktik mandiri. & Mahasiswa mengerjakan aktivitas terarah untuk persona dan journey map dan menunjukkan evidence hasil belajar. & Bentuk penilaian: Persona dan journey map. Mengacu rubrik RTM. & Ketepatan konsep; kualitas artefak/praktik; validasi dan dokumentasi. & 2.8 \\
4 & SCPMK0202-01201 & Kuis 1 konsep UI/UX. & Modalitas: Blended Learning. Bentuk: Luring/Daring. Metode: Design Thinking, studio review, studi kasus, dan peer feedback. & 1 x 2 x 50' tatap muka; 1 x 2 x 50' tugas/praktik mandiri. & Mahasiswa mengerjakan aktivitas terarah untuk kuis 1 konsep ui/ux dan menunjukkan evidence hasil belajar. & Bentuk penilaian: Kuis 1 konsep UI/UX. Mengacu rubrik RTM. & Ketepatan konsep; kualitas artefak/praktik; validasi dan dokumentasi. & 5 \\
5 & SCPMK0202-01202 & Figma dan komponen dasar. & Modalitas: Blended Learning. Bentuk: Luring/Daring. Metode: Design Thinking, studio review, studi kasus, dan peer feedback. & 1 x 2 x 50' tatap muka; 1 x 2 x 50' tugas/praktik mandiri. & Mahasiswa mengerjakan aktivitas terarah untuk figma dan komponen dasar dan menunjukkan evidence hasil belajar. & Bentuk penilaian: Figma dan komponen dasar. Mengacu rubrik RTM. & Ketepatan konsep; kualitas artefak/praktik; validasi dan dokumentasi. & 2.8 \\
6 & SCPMK0202-01202 & Wireframe dan information architecture. & Modalitas: Blended Learning. Bentuk: Luring/Daring. Metode: Design Thinking, studio review, studi kasus, dan peer feedback. & 1 x 2 x 50' tatap muka; 1 x 2 x 50' tugas/praktik mandiri. & Mahasiswa mengerjakan aktivitas terarah untuk wireframe dan information architecture dan menunjukkan evidence hasil belajar. & Bentuk penilaian: Wireframe dan information architecture. Mengacu rubrik RTM. & Ketepatan konsep; kualitas artefak/praktik; validasi dan dokumentasi. & 2.8 \\
7 & SCPMK0202-01202 & Visual design. & Modalitas: Blended Learning. Bentuk: Luring/Daring. Metode: Design Thinking, studio review, studi kasus, dan peer feedback. & 1 x 2 x 50' tatap muka; 1 x 2 x 50' tugas/praktik mandiri. & Mahasiswa mengerjakan aktivitas terarah untuk visual design dan menunjukkan evidence hasil belajar. & Bentuk penilaian: Visual design. Mengacu rubrik RTM. & Ketepatan konsep; kualitas artefak/praktik; validasi dan dokumentasi. & 2.8 \\
8 & SCPMK0202-01201, SCPMK0202-01202 & UTS artefak desain awal. & Modalitas: Blended Learning. Bentuk: Luring/Daring. Metode: Design Thinking, studio review, studi kasus, dan peer feedback. & 1 x 2 x 50' tatap muka; 1 x 2 x 50' tugas/praktik mandiri. & Mahasiswa mengerjakan aktivitas terarah untuk uts artefak desain awal dan menunjukkan evidence hasil belajar. & Bentuk penilaian: UTS artefak desain awal. Mengacu rubrik RTM. & Ketepatan konsep; kualitas artefak/praktik; validasi dan dokumentasi. & 20 \\
9 & SCPMK0202-01202 & Design system. & Modalitas: Blended Learning. Bentuk: Luring/Daring. Metode: Design Thinking, studio review, studi kasus, dan peer feedback. & 1 x 2 x 50' tatap muka; 1 x 2 x 50' tugas/praktik mandiri. & Mahasiswa mengerjakan aktivitas terarah untuk design system dan menunjukkan evidence hasil belajar. & Bentuk penilaian: Design system. Mengacu rubrik RTM. & Ketepatan konsep; kualitas artefak/praktik; validasi dan dokumentasi. & 2.8 \\
10 & SCPMK0202-01202 & UX writing. & Modalitas: Blended Learning. Bentuk: Luring/Daring. Metode: Design Thinking, studio review, studi kasus, dan peer feedback. & 1 x 2 x 50' tatap muka; 1 x 2 x 50' tugas/praktik mandiri. & Mahasiswa mengerjakan aktivitas terarah untuk ux writing dan menunjukkan evidence hasil belajar. & Bentuk penilaian: UX writing. Mengacu rubrik RTM. & Ketepatan konsep; kualitas artefak/praktik; validasi dan dokumentasi. & 2.8 \\
11 & SCPMK0601-01203 & Prototype interaktif. & Modalitas: Blended Learning. Bentuk: Luring/Daring. Metode: Design Thinking, studio review, studi kasus, dan peer feedback. & 1 x 2 x 50' tatap muka; 1 x 2 x 50' tugas/praktik mandiri. & Mahasiswa mengerjakan aktivitas terarah untuk prototype interaktif dan menunjukkan evidence hasil belajar. & Bentuk penilaian: Prototype interaktif. Mengacu rubrik RTM. & Ketepatan konsep; kualitas artefak/praktik; validasi dan dokumentasi. & 2.8 \\
12 & SCPMK0601-01203 & Kuis 2 design system dan prototype. & Modalitas: Blended Learning. Bentuk: Luring/Daring. Metode: Design Thinking, studio review, studi kasus, dan peer feedback. & 1 x 2 x 50' tatap muka; 1 x 2 x 50' tugas/praktik mandiri. & Mahasiswa mengerjakan aktivitas terarah untuk kuis 2 design system dan prototype dan menunjukkan evidence hasil belajar. & Bentuk penilaian: Kuis 2 design system dan prototype. Mengacu rubrik RTM. & Ketepatan konsep; kualitas artefak/praktik; validasi dan dokumentasi. & 5 \\
13 & SCPMK0601-01203 & UX motion. & Modalitas: Blended Learning. Bentuk: Luring/Daring. Metode: Design Thinking, studio review, studi kasus, dan peer feedback. & 1 x 2 x 50' tatap muka; 1 x 2 x 50' tugas/praktik mandiri. & Mahasiswa mengerjakan aktivitas terarah untuk ux motion dan menunjukkan evidence hasil belajar. & Bentuk penilaian: UX motion. Mengacu rubrik RTM. & Ketepatan konsep; kualitas artefak/praktik; validasi dan dokumentasi. & 2.8 \\
14 & SCPMK0601-01203 & Usability testing. & Modalitas: Blended Learning. Bentuk: Luring/Daring. Metode: Design Thinking, studio review, studi kasus, dan peer feedback. & 1 x 2 x 50' tatap muka; 1 x 2 x 50' tugas/praktik mandiri. & Mahasiswa mengerjakan aktivitas terarah untuk usability testing dan menunjukkan evidence hasil belajar. & Bentuk penilaian: Usability testing. Mengacu rubrik RTM. & Ketepatan konsep; kualitas artefak/praktik; validasi dan dokumentasi. & 2.8 \\
15 & SCPMK0802-01204 & Iterasi desain. & Modalitas: Blended Learning. Bentuk: Luring/Daring. Metode: Design Thinking, studio review, studi kasus, dan peer feedback. & 1 x 2 x 50' tatap muka; 1 x 2 x 50' tugas/praktik mandiri. & Mahasiswa mengerjakan aktivitas terarah untuk iterasi desain dan menunjukkan evidence hasil belajar. & Bentuk penilaian: Iterasi desain. Mengacu rubrik RTM. & Ketepatan konsep; kualitas artefak/praktik; validasi dan dokumentasi. & 2.8 \\
16 & SCPMK0202-01201--SCPMK0802-01204 & PBL proyek UI/UX. & Modalitas: Blended Learning. Bentuk: Luring/Daring. Metode: Design Thinking, studio review, studi kasus, dan peer feedback. & 1 x 2 x 50' tatap muka; 1 x 2 x 50' tugas/praktik mandiri. & Mahasiswa mengerjakan aktivitas terarah untuk pbl proyek ui/ux dan menunjukkan evidence hasil belajar. & Bentuk penilaian: PBL proyek UI/UX. Mengacu rubrik RTM. & Ketepatan konsep; kualitas artefak/praktik; validasi dan dokumentasi. & 11.4 \\
17 & SCPMK0601-01203, SCPMK0802-01204 & UAS validasi prototype. & Modalitas: Blended Learning. Bentuk: Luring/Daring. Metode: Design Thinking, studio review, studi kasus, dan peer feedback. & 1 x 2 x 50' tatap muka; 1 x 2 x 50' tugas/praktik mandiri. & Mahasiswa mengerjakan aktivitas terarah untuk uas validasi prototype dan menunjukkan evidence hasil belajar. & Bentuk penilaian: UAS validasi prototype. Mengacu rubrik RTM. & Ketepatan konsep; kualitas artefak/praktik; validasi dan dokumentasi. & 25 \\
\end{longtblr}
